% EDITION STATUS: partial; source units 19/58
\documentclass[11pt,a4paper]{article}
\usepackage[no-math]{fontspec}
\setmainfont{Latin Modern Roman}
\setsansfont{Latin Modern Sans}
\setmonofont{Latin Modern Mono}
\usepackage[french,english]{babel}
\usepackage[left=22mm,right=34mm,top=22mm,bottom=22mm,marginparwidth=23mm,marginparsep=5mm,includeheadfoot]{geometry}
\usepackage{amsmath,amssymb}
\usepackage[normalem]{ulem}
\usepackage{xcolor,tikz-cd,marginnote,etoolbox,microtype}
\usepackage[unicode,colorlinks=true,linkcolor=LinkInk,urlcolor=LinkInk,bookmarksopen=true]{hyperref}
\definecolor{LinkInk}{HTML}{254E70}
\definecolor{HandInk}{HTML}{254E70}
\definecolor{DeletionInk}{HTML}{8B3A3A}
\definecolor{QuietInk}{HTML}{555555}
\setlength{\parindent}{1em}
\setlength{\parskip}{3pt plus 1pt}
\setlength{\emergencystretch}{1.5em}
\raggedbottom
\normalmarginpar
\frenchspacing
\newif\ifdiplomatic
\newif\ifreading
\newif\ifenglish
\diplomaticfalse
\readingfalse
\englishtrue
\newcommand{\EditionTitle}{Sketch of a Program}
\newcommand{\EditionForm}{Corrected text with translator’s notes — annotated master}
\hypersetup{pdftitle={Sketch of a Program — Corrected text with translator’s notes},pdfauthor={Alexander Grothendieck},pdfsubject={A scan-checked edition of Esquisse d'un programme, January 1984, Montpellier cote 119}}
\newcommand{\DeclareScan}[4]{%
  \expandafter\def\csname edurl@#1\endcsname{#2}%
  \expandafter\def\csname edkey@#1\endcsname{#3}%
  \expandafter\def\csname edpage@#1\endcsname{#4}}
\def\CurrentScanURL{}
\def\CurrentPageKey{}
\def\CurrentCopy{copy1}
\def\CurrentApparatus{}
\newcounter{appnote}
\newcommand{\EditionPart}[2]{\clearpage\def\CurrentCopy{#1}\section*{#2}}
\newcommand{\scanpage}[1]{%
  \clearpage
  \expandafter\let\expandafter\CurrentScanURL\csname edurl@#1\endcsname
  \expandafter\let\expandafter\CurrentPageKey\csname edkey@#1\endcsname
  \def\CurrentApparatus{}%
  \hypertarget{scan-\CurrentPageKey}{}%
  \typeout{EDITION-SCAN: \CurrentPageKey; PDFPAGE: \the\numexpr\ReadonlyShipoutCounter+1\relax}%
  \ifreading
    \marginnote{\raggedright\scriptsize\sffamily\href{\CurrentScanURL}{119 / \csname edpage@#1\endcsname}}%
  \else
    \noindent{\small\sffamily\href{\CurrentScanURL}{[scan #1]}}\par\smallskip
  \fi}
\newcommand{\mspage}[1]{%
  \ifreading\marginnote{\raggedright\scriptsize\sffamily p.\,#1}[2\baselineskip]%
  \else\noindent{\small\sffamily[p.\,#1]}\par\smallskip\fi}
\newcommand{\hand}[1]{\ifdiplomatic\textcolor{HandInk}{#1}\else#1\fi}
\let\over\relax
\newcommand{\over}[2]{\textcolor{QuietInk}{#1}\textcolor{HandInk}{\ensuremath{^{\!\nearrow}}#2}}
\newcommand{\inserted}[2][]{\textcolor{HandInk}{\ensuremath{\langle}#2\ensuremath{\rangle}}}
\newcommand{\deleted}[1]{\ifmmode\textcolor{DeletionInk}{[#1]_{\scriptscriptstyle\mathrm{del}}}\else\textcolor{DeletionInk}{\sout{#1}}\fi}
\renewcommand{\xout}[1]{\ifmmode\textcolor{QuietInk}{[#1]_{\scriptscriptstyle\mathrm{x}}}\else\textcolor{QuietInk}{\sout{#1}}\fi}
\let\TextUline\uline
\let\TextUuline\uuline
\renewcommand{\uline}[1]{\ifmmode\underline{#1}\else\TextUline{#1}\fi}
\renewcommand{\uuline}[1]{\ifmmode\underline{\underline{#1}}\else\TextUuline{#1}\fi}
\newcommand{\unsure}[1]{\href{\CurrentScanURL}{[#1?]}}
\newcommand{\illegible}[1]{\href{\CurrentScanURL}{[\ifenglish illegible\else illisible\fi\ #1]}}
\newsavebox{\MarginalTestBox}
\newcommand{\margin}[1]{%
  \sbox{\MarginalTestBox}{\parbox{\marginparwidth}{\raggedright\footnotesize#1}}%
  \ifdim\dimexpr\ht\MarginalTestBox+\dp\MarginalTestBox\relax>3\baselineskip
    \par\smallskip\noindent\begin{minipage}{\linewidth}\hrule\smallskip
    {\footnotesize\textcolor{HandInk}{\ifenglish Marginal note:\else Note marginale :\fi\ #1}}\smallskip\hrule\end{minipage}\par\smallskip
  \else
    \marginnote{\raggedright\footnotesize\textcolor{HandInk}{#1}}%
  \fi}
\newcommand{\authornote}[1]{\margin{#1}}
\newcommand{\diagram}[1]{\par\smallskip{\centering#1\par}\smallskip}
\newenvironment{diagramblock}{\par\begin{center}}{\end{center}\par}
\newcommand{\DeclareVariant}[2]{\expandafter\def\csname edvariant@#1\endcsname{#2}}
\newcommand{\AppItem}[2]{\par\noindent\textsuperscript{#1}\csname edvariant@#2\endcsname}
\newcommand{\EditionApp}[1]{%
  \stepcounter{appnote}%
  \ifmmode{}^{\scriptscriptstyle\theappnote}\else\textsuperscript{\theappnote}\fi
  \begingroup\edef\x{\endgroup\noexpand\gappto\noexpand\CurrentApparatus{\noexpand\AppItem{\theappnote}{#1}}}\x}
\newcommand{\var}[4]{\ifdiplomatic#1\else#2\ifreading\else\EditionApp{#4}\fi\fi}
\newcommand{\copytwo}[2]{\ifdiplomatic\else\ifreading\else\DeclareVariant{#1}{#2}\EditionApp{#1}\fi\fi}
\newcommand{\tnote}[2]{\ifreading\else\DeclareVariant{#1}{\textit{Translator's note.} #2}\EditionApp{#1}\fi}
\newcommand{\finishscan}{\ifx\CurrentApparatus\empty\else\par\bigskip\noindent\rule{.35\linewidth}{.35pt}\par\smallskip{\footnotesize\setlength{\parskip}{3pt}\CurrentApparatus}\fi}
\newcommand{\DeclareNoteReturns}[2]{\expandafter\def\csname edreturns@#1\endcsname{#2}}
\newcommand{\notecall}[3]{\hypertarget{#2}{\hyperlink{note-#1}{#3}}}
\newcommand{\notedef}[2]{\hypertarget{note-#1}{#2}\,{\scriptsize[\csname edreturns@#1\endcsname]}}
\newcommand{\AppReading}[1]{{\def\margin##1{##1}\def\authornote##1{##1}#1}}
\newcommand{\EmptyReading}{\ifenglish\textit{omitted}\else\textit{supprimé}\fi}
\DeclareScan{119/2}{https://grothendieck.umontpellier.fr/web/viewer.html?file=../119.pdf\#page=3}{esq_119_002}{2}
\DeclareScan{119/3}{https://grothendieck.umontpellier.fr/web/viewer.html?file=../119.pdf\#page=4}{esq_119_003}{3}
\DeclareScan{119/4}{https://grothendieck.umontpellier.fr/web/viewer.html?file=../119.pdf\#page=5}{esq_119_004}{4}
\DeclareScan{119/5}{https://grothendieck.umontpellier.fr/web/viewer.html?file=../119.pdf\#page=6}{esq_119_005}{5}
\DeclareScan{119/6}{https://grothendieck.umontpellier.fr/web/viewer.html?file=../119.pdf\#page=7}{esq_119_006}{6}
\DeclareScan{119/7}{https://grothendieck.umontpellier.fr/web/viewer.html?file=../119.pdf\#page=8}{esq_119_007}{7}
\DeclareScan{119/8}{https://grothendieck.umontpellier.fr/web/viewer.html?file=../119.pdf\#page=9}{esq_119_008}{8}
\DeclareScan{119/9}{https://grothendieck.umontpellier.fr/web/viewer.html?file=../119.pdf\#page=10}{esq_119_009}{9}
\DeclareScan{119/10}{https://grothendieck.umontpellier.fr/web/viewer.html?file=../119.pdf\#page=11}{esq_119_010}{10}
\DeclareScan{119/11}{https://grothendieck.umontpellier.fr/web/viewer.html?file=../119.pdf\#page=12}{esq_119_011}{11}
\DeclareScan{119/12}{https://grothendieck.umontpellier.fr/web/viewer.html?file=../119.pdf\#page=13}{esq_119_012}{12}
\DeclareScan{119/13}{https://grothendieck.umontpellier.fr/web/viewer.html?file=../119.pdf\#page=14}{esq_119_013}{13}
\DeclareScan{119/14}{https://grothendieck.umontpellier.fr/web/viewer.html?file=../119.pdf\#page=15}{esq_119_014}{14}
\DeclareScan{119/15}{https://grothendieck.umontpellier.fr/web/viewer.html?file=../119.pdf\#page=16}{esq_119_015}{15}
\DeclareScan{119/16}{https://grothendieck.umontpellier.fr/web/viewer.html?file=../119.pdf\#page=17}{esq_119_016}{16}
\DeclareScan{119/17}{https://grothendieck.umontpellier.fr/web/viewer.html?file=../119.pdf\#page=18}{esq_119_017}{17}
\DeclareScan{119/18}{https://grothendieck.umontpellier.fr/web/viewer.html?file=../119.pdf\#page=19}{esq_119_018}{18}
\DeclareScan{119/19}{https://grothendieck.umontpellier.fr/web/viewer.html?file=../119.pdf\#page=20}{esq_119_019}{19}
\DeclareScan{119/20}{https://grothendieck.umontpellier.fr/web/viewer.html?file=../119.pdf\#page=21}{esq_119_020}{20}
\DeclareVariant{V-FR-010-001}{\textit{slip corrected}. \AppReading{Ladegraillerie} \(\rightarrow\) \AppReading{Ladegaillerie}. The preceding line gives Yves Ladegaillerie for the same participant; the second occurrence interposes an r.}

\begin{document}
\begin{titlepage}
\centering
\vspace*{18mm}
{\Large Alexander Grothendieck\par}
\vspace{12mm}
{\Huge\bfseries\EditionTitle\par}
\vspace{8mm}
{\Large\EditionForm\par}
\vspace{14mm}
{\large January 1984\par}
\vfill
{\large Université de Montpellier\par fonds Grothendieck, cote 119\par}
\vspace{8mm}
{\normalsize Two typescript copies: sites 2--59 and 78--135\par}
\vspace{8mm}
{\normalsize Edition dated 14 November 2026\par}
\end{titlepage}
\thispagestyle{plain}
\section*{Preface}
\noindent\textbf{Partial working edition.} This file contains 19 of the 58 source pages required for this form. Collation and final validation remain in progress.\par\smallskip
This English translation follows the corrected French text of copy 1. Authorial changes and justified corrections of isolated typing slips are recorded in the apparatus. Copy 2 annotations are identified separately. Translator’s notes explain words, wordplay, abbreviations and notation at the place where an explanation is needed. Unresolved readings remain marked and linked to their scans. The author’s notes (1)–(7) remain at the end, with links in both directions. The formulas retain the exact French source strings; terminology follows the supplied unified English glossary.

\smallskip
AI-transcribed by Claude Opus 5.5 (Anthropic, effort high, Message Batches API); checked against the scans, and this edition built, by GPT-6 Astra (OpenAI, ChatGPT Pro mode); prompts and records by Claude (Anthropic).
AI-translated by GPT-6 Astra (OpenAI, ChatGPT Pro mode).

The transcription of copy 1 was made with the text layer of the Grothendieck Circle's typeset \textit{Esquisse} (\texttt{EsquisseFr.pdf}) given to the transcribing model as a draft, to compare with after it had transcribed each page; every reading of this edition was decided on the scans.

The published French text and English translation in \textit{Geometric Galois Actions 1}, edited by L. Schneps and P. Lochak (Cambridge University Press, 1997), the Grothendieck Circle's typeset \texttt{EsquisseFr.pdf}, and the Centre for Grothendieckian Studies' draft transcription G94d are named as published texts. They are not reproduced here.

\section*{Rights and credits}
The rights in the original work belong to its author or the author's representatives or heirs. Where we hold rights, in the AI-produced transcription, translation, corrections and typesetting, they are dedicated to the public domain under CC0 1.0. This edition is published for educational, non-profit purposes, and we claim no rights we do not hold. If you hold rights in this work and object to its publication here, please open a GitHub issue and we will take it down.

\clearpage

% BEGIN PAGE esq_119_002
\scanpage{119/2}
\uline{SKETCH OF A PROGRAM}

by Alexandre Grothendieck

\uline{Contents} :

1. Opening.\tnote{TN-002-01}{The French heading is ``Envoi'', a word for a sending or dispatch, used here to introduce the text addressed to the application committee and other readers.}

2. A "Légo-Teichmüller game", and the Galois group of $\overline{\mathbb{Q}}$ over $\mathbb{Q}$.

3. Number fields associated with a child's drawing.\tnote{TN-002-02}{``Child's drawing'' translates ``dessin d'enfant''. The phrase becomes a mathematical term in the discussion.}

4. Regular polyhedra and finite fields.

5. A hue and cry against so-called ``general'' topology, and heuristic reflections
toward a so-called ``tame'' topology.\tnote{TN-002-03}{``Haro sur'' is a denunciatory cry directed against something; ``modérée'' is rendered ``tame'' throughout, following the terminology standard supplied for this edition.}

6. ``Differentiable theories'' (in the manner of Nash) and ``tame theories''.

7. Pursuing Stacks.\tnote{TN-002-04}{The French title is ``A la Poursuite des Champs''. ``Champ'' is translated ``stack'' in accordance with the supplied terminology standard.}

8. Digressions in two-dimensional geometry.

9. An appraisal of teaching activity.

10. Epilogue.

Notes.

N.B. The asterisks (*) refer to the notes appearing at the bottom of
the same page, the numbered references from (\textsuperscript{\hand{1}}) to (\textsuperscript{\hand{7}}\unsure{.}) to the notes (added
later) collected at the end of the report.

\mspage{\hand{2}}

\finishscan
% END PAGE esq_119_002
% BEGIN PAGE esq_119_003
\scanpage{119/3}
\mspage{1}

SKETCH OF A PROGRAM

by A. Grothendieck

1. As the present circumstances make the prospects of teaching at research level
at the University increasingly illusory for me, I have resolved to apply for
admission to the CNRS,\tnote{TN-003-01}{CNRS: \href{https://aissai.cnrs.fr/en/legal-notice/}{Centre national de la recherche scientifique}, the French National Center for Scientific Research.}
so that I may devote my energy to developing work and perspectives
that, it is becoming clear, no student (nor even, it seems,
any fellow mathematician) will be found to develop in my
stead.

In place of a ``Qualifications and Works'' document,\tnote{TN-003-02}{``Titres et Travaux'' is the French designation of the document listing qualifications and scholarly work.}
following this text will be found the full reproduction of a sketch, arranged by themes,
of what I regarded as my main mathematical contributions at the time
of writing that report, in 1972. It also contains
a list of articles published by that date. I have published no scientific
articles since 1970. In the lines that follow,
I propose to give at least a glimpse of some principal themes
of my mathematical reflections since then. Over the years these reflections
have taken material form in two bulky boxes of
manuscript notes, probably difficult for anyone other than
myself to decipher, which, after successive stages of decantation,
are perhaps awaiting their moment for an overall write-up, at least
a provisional one, intended for the mathematical community. The
term ``write-up'' here is somewhat inappropriate,\tnote{TN-003-03}{``Write-up'' renders ``rédaction''. The source itself immediately explains the continuing work of discovery meant by this word.}
since it is much more a matter of developing the many ideas and visions begun
over these last twelve years, making them more precise and exploring them
more deeply, with all the unforeseen turns that constantly
accompany this kind of work - a work of discovery, then, and
not of compiling notes piously accumulated. And I fully intend,
in writing the ``Mathematical Reflections'' begun in February
1983,\tnote{TN-003-04}{The French title is ``Réflexions Mathématiques''.}
to let the course of thought that probes and discovers appear
clearly through the pages, often groping in the half-light,
with sudden openings of light when some stubborn

\margin{3}

\finishscan
% END PAGE esq_119_003
% BEGIN PAGE esq_119_004
\scanpage{119/4}
\mspage{2}

false, or merely inadequate, image is finally flushed out
and brought to light, and the things that seemed askew
fall into place, in the mutual harmony proper to them.
\qquad Be that as it may, the following sketch of some
themes of reflection from the last ten or twelve years will
serve at the same time as a sketch of a work program for the
coming years, which I intend to devote to developing
these themes, or at least some of them. It is intended,
on the one hand, for the colleagues on the National Committee called upon to rule
on my application,\tnote{TN-004-01}{``National Committee'' translates ``Comité National'', the committee named in the source in connection with the CNRS application.}
and, on the other, for a few other colleagues, former
students, and friends, in the event that some of the ideas sketched
here might interest one of them.

2. \qquad The demands of university teaching, addressed
therefore to students (including so-called ``advanced'' students) whose
mathematical background is modest (and often less than modest), have led me
to renew drastically the themes of reflection to propose
to my students, and, one thing leading to another and increasingly, to myself
as well. It had seemed important to me to start from a common
intuitive background, independent of any technical language meant to express it,
far earlier than any such language - it turned out that the geometric
and topological intuition of shapes, and more particularly of
two-dimensional shapes, was such a common ground. Thus the themes in question
can be grouped under the name ``topology of surfaces'' or
``geometry of surfaces'', with the understanding, in the latter name,
that the principal emphasis falls on the topological properties
of surfaces, or on the combinatorial aspects that constitute
their most down-to-earth technical expression, and not on the differential
aspects, or even the conformal, riemannian, holomorphic and (thence)
the ``complex algebraic curves'' aspect. Once this last step
has been taken, however, suddenly algebraic geometry (my old
love !) bursts in again, through the very objects that
can be regarded as the ultimate building blocks of all
the other algebraic varieties. Whereas in my research before

\margin{4}

\finishscan
% END PAGE esq_119_004
% BEGIN PAGE esq_119_005
\scanpage{119/5}
\mspage{3}

1970, my attention was systematically directed toward objects
of the greatest generality, in order to draw out an adequate overall language
for the world of algebraic geometry, and I lingered
over algebraic curves only to the strict extent that this proved
indispensable (particularly in étale cohomology) for developing
``all-purpose'' techniques and statements valid in every dimension
and in every place (I mean, over all base schemes, or even
all ringed base topoi ...), here I am thus brought back, through
objects so simple that a child can come to know them through play,
to the beginnings and origins of algebraic geometry, familiar to Riemann
and his followers !

Since about 1975, then, the geometry of (real)
surfaces, and from 1977 onward the connections between questions in
the geometry of surfaces and the algebraic geometry of algebraic curves
defined over fields such as $\mathbb{C}$, $\mathbb{R}$ or extensions of finite
type of $\mathbb{Q}$, have been my principal source of inspiration, as well
as my constant guiding thread. It was with surprise and wonder
that over the years I discovered (or rather, probably,
rediscovered) the prodigious, truly inexhaustible richness, the
unsuspected depth of this seemingly so innocuous theme. I believe
I sense in it a particularly vital point, a privileged point of convergence
of the main currents of mathematical ideas, as also
of the main structures and the visions of things they express,
from the most specific (such as the rings $\mathbb{Z}$, $\mathbb{Q}$, $\bar{\mathbb{Q}}$, $\mathbb{R}$, $\mathbb{C}$,
or the group $S\ell(2)$ over one of these rings, or general reductive
algebraic groups) to the most ``abstract'', such as algebraic,
complex analytic or real analytic ``multiplicities''.\tnote{TN-005-01}{``Multiplicity'' consistently renders the author's ``multiplicité''. His parenthesis here explains its role in retaining the automorphisms of the objects under study.}
(These enter naturally when one sets out to study
``moduli varieties'' systematically for the geometric objects
under consideration, seeking to go beyond the notoriously insufficient
viewpoint of ``coarse moduli'', which amounts, most unfortunately,
to ``killing'' the automorphism groups of these objects .) Among these
moduli multiplicities, it is those of Mumford-Deligne for
``stable'' algebraic curves of genus $g$, with $\nu$ marked points, which

\margin{5}

\finishscan
% END PAGE esq_119_005
% BEGIN PAGE esq_119_006
\scanpage{119/6}
\mspage{4}

I denote by $M^{\wedge}_{g,\nu}$ (the compactification of the ``open'' multiplicity
$M_{g,\nu}$ corresponding to smooth curves), that for some two
or three years have exerted a particular fascination upon me,
perhaps stronger than that of any other mathematical object to date.
To tell the truth, it is rather the system of \uline{all} the multiplicities
$M_{g,\nu}$ for varying $g,\nu$, connected with one another by a certain number
of fundamental operations (such as the operations of ``filling
holes'', i.e. ``erasing'' marked points, that of ``gluing'',
and the inverse operations), which reflect, in absolute algebraic
geometry of characteristic zero (for the moment), familiar geometric
operations from the viewpoint of topological or conformal ``surgery''
on surfaces. The principal reason for this fascination is probably
that this very rich geometric structure on the
system of ``open'' moduli multiplicities $M_{g,\nu}$ is reflected
in an analogous structure on the corresponding fundamental
groupoids, the ``Teichmüller groupoids'' $\hat{T}_{g,\nu}$ ,and that these operations
at the level of the $\hat{T}_{g,\nu}$ have a sufficiently intrinsic character
for the Galois group $\mathrm{I}\!\Gamma$ of $\overline{\mathbb{Q}}/\mathbb{Q}$ to act on this entire ``tower'' of
Teichmüller groupoids, respecting all these structures. More
extraordinary still, this action is \uline{faithful} - in fact,
it is already faithful on the first nontrivial ``level'' of this tower,
namely $\hat{T}_{0,4}$ - which also means, essentially, that the outer
action of $\mathrm{I}\!\Gamma$ on the fundamental group $\hat{\pi}_{0,3}$ of the standard
projective line $\mathbb{P}^1$ over $\mathbb{Q}$, with the three points $0,\ 1,\ \infty$ removed, is
already faithful. Thus \uline{the Galois group} $\mathrm{I}\!\Gamma$ \uline{is realized as a group}
\uline{of automorphisms of a most concrete profinite group}, moreover respecting
certain essential structures of that group. It follows
that an element of $\mathrm{I}\!\Gamma$ can be ``parametrized'' (in various equivalent
ways, moreover) by a suitable element of this profinite group
$\hat{\pi}_{0,3}$ (a free profinite group on two generators), or by a system
of such elements, this element or these elements being, moreover, subject
to certain simple conditions, necessary (and probably insufficient)
for this element or these elements actually to correspond to an element of $\mathrm{I}\!\Gamma$.
One of the most fascinating tasks here is precisely to grasp
necessary \uline{and} sufficient conditions on an outer automorphism

\margin{6}

\finishscan
% END PAGE esq_119_006
% BEGIN PAGE esq_119_007
\scanpage{119/7}
\mspage{5}

of $\hat{\pi}_{0,3}$, i.e. on the corresponding parameter or parameters, for it
to come from an element of $\mathrm{I}\!\Gamma$ - which would provide a
``purely algebraic'' description, in terms of profinite groups and without
reference to the Galois theory of number fields, of the Galois
group $\mathrm{I}\!\Gamma = \mathrm{Gal}(\bar{\mathbb{Q}}/\mathbb{Q})$ !

Perhaps even a conjectural characterization of $\mathrm{I}\!\Gamma$
as a subgroup of Autext($\hat{\pi}_{0,3}$) is for the moment out of
reach\notecall{copy1-1}{call-esq_119_007-1-1}{($^1$)} ; I have no conjecture to propose yet. Another
task, however, can be approached immediately: describing
the action of $\mathrm{I}\!\Gamma$ on the entire Teichmüller tower in terms of its
action on the ``first level'' $\hat{\pi}_{0,3}$ , i.e. expressing an automorphism
of this tower in terms of the ``parameter'' in $\hat{\pi}_{0,3}$ that specifies
the varying element $\gamma$ of $\mathrm{I}\!\Gamma$. This is connected with a representation of the
Teichmüller tower (as a groupoid equipped with a ``gluing''
operation) by generators and relations, which will in particular give
a presentation by generators and relations, in the ordinary sense, of
each of the $\hat{T}_{g,\nu}$ (as a profinite groupoid). Here, even for
$g = 0$ (thus when the corresponding Teichmüller groups are
``well-known'' braid groups), the generators and relations
known to date that have come to my attention seem unusable
as they stand, because they lack the invariance and
symmetry properties indispensable for the action of $\mathrm{I}\!\Gamma$ to be directly
readable from this presentation. This is connected in particular with the fact that
people still persist, when calculating with fundamental
groups, in fixing a single base point, instead of cleverly choosing
a whole collection invariant under the symmetries of the
situation, which are thereby lost along the way. In certain situations
(such as descent theorems in the manner of Van Kampen for fundamental
groups), it is much more elegant, even indispensable for understanding
anything about them, to work with fundamental groupoids
relative to a suitable collection of base points , and this is certainly
so for the Teichmüller tower. It would seem (incredible,
but true !) that the very geometry of the first level of the
Teichmüller tower (thus corresponding to ``moduli'' either for projective

\margin{7}

\finishscan
% END PAGE esq_119_007
% BEGIN PAGE esq_119_008
\scanpage{119/8}
\mspage{6}

lines with four marked points or for elliptic
curves (!)) has never been properly made explicit, for example the relation
between the genus 0 case and the geometry of the octahedron and of the
tetrahedron. A fortiori, the moduli multiplicities $M_{0,5}$ (for
projective lines with five marked points) and $M_{1,2}$ (for
curves of genus 1 with two marked points), which are moreover almost
isomorphic to one another, seem virgin territory - braid groups
will not enlighten us about them ! I have begun to
look at $M_{0,5}$ in spare moments; it is a true jewel, with a
very rich geometry closely connected with that of the icosahedron.

The a priori interest of a complete knowledge of the first two
levels of the tower (namely, the cases in which the moduli dimension
$N = 3g-3+\nu$ is $\leq 2$) lies in this principle, that \uline{the entire tower
is reconstructed from the first two levels}, in the sense that, through
the fundamental ``gluing'' operation, level 1 supplies a
complete system of generators, and level 2 a complete system of
relations. There is a striking analogy, and I am convinced it is
more than formal, between this principle and Demazure's analogous
principle for the structure of reductive algebraic groups, when
the term ``level'' or ``moduli dimension'' is replaced by ``semisimple
rank of the reductive group''. The connection becomes still more striking
when one recalls that the Teichmüller group $T_{1,1}$ (now in the
discrete transcendental context, and not in the profinite algebraic
context, where one finds the profinite completions of the former) is
precisely $S\ell(2,\mathbb{Z})$, i.e. the group of integral points of the
``absolute'' simple group scheme of rank 1 $S\ell(2)_{\mathbb{Z}}$. Thus, \uline{the fundamental
building block for the Teichmüller tower} \uline{is essentially
the same as that for ``the tower'' of reductive groups of all ranks} -
a group of which one can probably say, moreover, that it is present
in all the essential disciplines of mathematics.

This principle for constructing the Teichmüller tower has
not been proved at present - but I have no doubt that it
is valid. It would follow (through a theory of dévissage of stratified
structures - in this instance the $M^{\wedge}_{g,\nu}$ - still to be written,
cf. par. 5) from an extremely plausible property of the open

\margin{8}

\finishscan
% END PAGE esq_119_008
% BEGIN PAGE esq_119_009
\scanpage{119/9}
\mspage{7}

\unsure{moduli}\tnote{TN-009-01}{The checked French retains the alternatives ``mcdulaires|modulaires'' for the unclear typed letter. Both lead to ``moduli'' in this phrase; the doubt remains marked and linked to the scan.}
multiplicities $M_{g,\nu}$ in the complex analytic context,
namely that, for a moduli dimension $N \geq 3$, the fundamental
group of $M_{g,\nu}$ (i.e. the usual Teichmüller group $T_{g,\nu}$)
is isomorphic to the ``fundamental group at infinity'', i.e. that of a
``tubular neighborhood of infinity''. This is a very familiar fact
(due essentially to Lefschetz) for a smooth \uline{affine} variety of
dimension $N \geq 3$. It is true that the moduli multiplicities are
not affine (except for small values of $g$), but it would suffice
for such an $M_{g,\nu}$ of dimension $N$ (or rather, a suitable finite
covering) to be the union of $N$-2 affine open sets, thus
for $M_{g,\nu}$ not to be ``too close to a compact variety''.

Having no doubt about this principle for constructing the
Teichmüller tower, I prefer to leave it to the experts in the transcendental
theory, better equipped than I am, to prove what is needed
(should any be interested), and instead to make explicit,
with all the care it deserves, the resulting structure of the
Teichmüller tower by generators and relations, this time in the
discrete setting and not the profinite one - which amounts, essentially,
to a complete understanding of the four moduli multiplicities
$M_{0,4}$, $M_{1,1}$, $M_{0,5}$, $M_{1,2}$, and of their fundamental groupoids relative
to suitably chosen ``base points''. These present themselves
quite naturally as the complex algebraic curves of the type
$(g,\nu)$ under consideration whose automorphism group (necessarily
finite) is larger than in the generic case (*). Including among them the
holomorphic sphere with three marked points (coming from $M_{0,3}$, i.e.

\rule{4cm}{0.4pt}

(*) One must also add the ``base points'' arising from operations
of gluing ``pieces'' of the same type in lower moduli dimension.
On the other hand, in moduli dimension 2 (the cases of $M_{0,5}$ and $M_{1,2}$),
one should exclude the points of certain one-parameter families
of curves admitting an exceptional automorphism of order 2. These
families, moreover, constitute remarkable rational curves on the multiplicities
under consideration, which seem to me an important ingredient
of the structure of these multiplicities.

\margin{9}

\finishscan
% END PAGE esq_119_009
% BEGIN PAGE esq_119_010
\scanpage{119/10}
\mspage{8}

from level O), one finds \uline{twelve} \uline{fundamental} ``\uline{building blocks}''
(6 of genus O, 6 of genus 1) in a ``Légo-Teichmüller game (large
box), where the marked points on the surfaces under consideration are replaced
by ``holes'' with boundary, so as to obtain surfaces with boundary,
thus building blocks that can be fitted together with gentle
friction, as in the ordinary Légo game dear to our children
(or grandchildren...). Assembling them provides a thoroughly
visual way of constructing every type of surface (it is essentially
these assemblies that will be the ``base points'' for
our famous tower), and also of visualizing the \uline{elementary} ``\uline{paths}''
through equally concrete operations, such as ``twists'', or
automorphisms of the pieces in the game, and of writing the \uline{fundamental relations}
between composite paths. Depending on the size (and the price !) of the
construction box used, one moreover \unsure{finds}\tnote{TN-010-01}{The French reading is retained as ``trcuve|trouve'' because the letterform is unresolved. Both alternatives are represented here by ``finds'', with the source doubt retained.}
many different descriptions of the Teichmüller tower by generators and
relations. The smallest box is reduced to pieces that are all
identical, of type (O,3) - these are Thurston's ``pairs of pants'',\tnote{TN-010-02}{``Pairs of pants'' translates ``pantalons'', the surface name used by the author. The letter O in the scanned type expressions is retained.}
and the Légo-Teichmüller game I am trying to describe, arising from
motivations and reflections in absolute algebraic geometry over the field $\mathbb{Q}$,
is very close to Thurston's game of ``hyperbolic geodesic surgery'',
whose existence I learned of last year from Yves Ladegaillerie.
In a microseminar with Carlos Contou-Carrère and Yves \var{Ladegraillerie}{Ladegaillerie}{slip corrected}{V-FR-010-001},
we have begun a reflection one of whose aims is to confront
the two viewpoints, which complement each other.

I add that each of the twelve building blocks in
the ``large box'' is equipped with a canonical cellular decomposition,
stable under all symmetries, whose only vertices are the
``marked points'' (or centers of the holes), and whose edges are certain
geodesic paths (for the canonical riemannian structure on the sphere
or torus under consideration) between certain pairs of vertices (namely those
that lie on the same ``real locus'', for a suitable real structure
on the complex algebraic curve under consideration). Consequently,
all the surfaces obtained in this game by assembly are equipped
with canonical cellular structures, which in their turn (cf.§3 below)

\margin{10}

\finishscan
% END PAGE esq_119_010
% BEGIN PAGE esq_119_011
\scanpage{119/11}
\mspage{9}

allow these surfaces to be considered as associated with
canonically determined complex algebraic curves (and even curves over $\overline{\mathbb{Q}}$).
Here there is a typical interplay between the combinatorial
and the complex algebraic (or better, the algebraic over $\overline{\mathbb{Q}}$).

The ``small box'', with its identical pieces, which has
the charm of economy, will probably give a relatively
complicated description of the relations (complicated, but by no means
inextricable !). The large box will give rise to more numerous
relations (because there are many more base points and
remarkable paths between them), but with a more transparent structure. I
expect that in moduli dimension 2, just as in the more or
less familiar case of moduli dimension 1 (with, in particular, the
description of $S\ell(2,\mathbb{Z})$ by $(\rho,\sigma \mid \rho^3 = \sigma^2,\ \rho^4 = \sigma^6 = 1)$), one will find
generation by the automorphism groups of the three relevant types
of pieces, with simple relations that I have not worked out
at the time of writing these lines. Perhaps one will even find a
principle of this kind for all the $T_{g,\nu}$, as well as a cellular
decomposition of $M^{\wedge}_{g,\nu}$ generalizing those that present themselves spontaneously
for $M^{\wedge}_{0,4}$ and $M^{\wedge}_{1,1}$, and that I already glimpse for
moduli dimension 2, using the hypersurfaces corresponding
to the various \uline{real structures} on the complex structures
under consideration to carry out the desired cellular decomposition.


3. Instead of following (as planned) a rigorous thematic
order, I have let myself be carried away by my predilection for a
particularly rich and burning theme, to which I moreover intend to
devote myself as a priority for some time, beginning with the start
of the 84/85 academic year. I therefore resume the thematic account where I left it,
at the very beginning of the preceding paragraph.

My interest in topological surfaces begins to emerge
in 1974, when I propose to Yves Ladegaillerie the theme of studying
isotopically the embeddings of a topological 1-complex in a compact
surface. In the two years that follow, this study leads him to
a remarkable isotopy theorem, giving a complete algebraic description

\margin{11}

\finishscan
% END PAGE esq_119_011
% BEGIN PAGE esq_119_012
\scanpage{119/12}
\mspage{10}

of the isotopy classes of embeddings of such 1-complexes,
or of compact surfaces with boundary, in a compact oriented surface,
in terms of certain very simple combinatorial invariants and the
fundamental groups of the participants. This theorem, which should
extend without difficulty to embeddings of any compact space
(triangulable, to simplify matters) in a compact oriented surface,
recovers as easy corollaries several deep classical results
of surface theory, and in particular Baer's isotopy theorem.
It is finally going to be published, separately from the rest (and
ten years later, given the harshness of the times...), in Topology. Ladegaillerie's
work also contains a purely algebraic description,
in terms of fundamental groupoids, of the
``isotopy'' category of compact surfaces $X$, equipped with a topological
1-complex $K$ embedded in $X$. This description, which had the misfortune
to go against the ``taste of the day'' and consequently seems unpublishable,
has nevertheless served (and still serves) as a valuable guide in
my subsequent reflections, particularly in the context of absolute
algebraic geometry of characteristic zero.

The case in which $(X,K)$ is a 2-dimensional ``map'',\tnote{TN-012-01}{``Map'' here translates ``carte'', a surface equipped with the embedded graph and cell conditions specified in the sentence.}
i.e. in which the connected components of $X\backslash K$ are open 2-cells \unsure{.}
(and in which $K$ is moreover equipped with a finite set $S$ of ``vertices'',
such that the connected components of $K\backslash S$ are open 1-cells)
gradually attracts my attention in the following years.
The isotopy category of these maps admits a particularly
simple algebraic description, through the set of ``frames'' (or ``flags'',
or ``biarcs'') associated with the map, which is naturally equipped
with the structure of a set with a group of operators, under the group
\[ \underline{C}_2 = \langle \sigma_0,\sigma_1,\sigma_2 \,|\, \sigma_0^{\,2} = \sigma_1^{\,2} = \sigma_2^{\,2} = (\sigma_0\sigma_2)^2 = 1 \rangle \ , \]
which I call the (unoriented) \uline{cartographic group} of dimension 2.
It admits as a subgroup of index 2 the \uline{oriented cartographic group},
generated by products of even numbers of the generators, which can
also be described as
\[ \underline{C}_2^{+} = \langle \rho_s,\rho_f,\sigma \,|\, \rho_s\rho_f = \sigma,\ \sigma^2 = 1 \rangle \ , \]

\margin{12}

\finishscan
% END PAGE esq_119_012
% BEGIN PAGE esq_119_013
\scanpage{119/13}
\mspage{11}

(with
\[ \rho_s = \sigma_2\sigma_1 \ , \quad \rho_f = \sigma_1\sigma_0 \ , \quad \sigma = \sigma_0\sigma_2 = \sigma_2\sigma_0 \ , \]
the operations of \uline{elementary rotation} of a frame about a vertex,
a face and an edge respectively). There is a perfect dictionary
between the topological situation of compact maps, resp.
compact oriented maps, on the one hand, and finite sets with
operator group $\underline{C}_2$ resp. $\underline{C}_2^+$ on the other, a dictionary whose
existence was, moreover, more or less known, but never
stated with the necessary precision or developed to any extent.
This foundational work is carried out with the care it deserves in
an excellent DEA project,\tnote{TN-013-01}{DEA abbreviates ``diplôme d'études approfondies'', an advanced studies diploma.}
undertaken jointly by Jean Malgoire and
Christine Voisin in 1976.

This reflection suddenly takes on a new dimension
with the simple observation that the group $\underline{C}_2^+$ can be interpreted
as a quotient of the fundamental group of an oriented sphere with
three points removed, numbered 0, 1, 2, the operations $\rho_s$, $\sigma$, $\rho_f$
being interpreted as the loops about these points, satisfying
the familiar relation
\[ \ell_0\ell_1\ell_2 = 1 \ , \]
while the additional relation $\sigma^2 = 1$, i.e. $\ell_1^{\,2} = 1$, means
that one is interested in the quotient of the fundamental group corresponding
to an imposed ramification index 2 at the point 1, which therefore classifies
coverings of the sphere ramified at most at the
points 0, 1, 2, with ramification equal to 1 or 2 at the points
above 1. Thus compact oriented maps form an
isotopy category equivalent to that of these coverings, subject
to the further condition of being finite coverings.
Taking now as the reference sphere the Riemann sphere,
or complex projective line, rigidified by the three points 0, 1
and $\infty$ (the last thus replacing 2), and recalling that every
finite ramified covering of a complex algebraic curve itself inherits
the structure of a complex algebraic curve, one arrives at this
observation, which eight years later still seems to me just as
extraordinary: \uline{every oriented} ``\uline{finite}'' \uline{map is canonically realized}
\uline{on a complex algebraic curve} ! Better still, since the complex projective

\margin{13}

\finishscan
% END PAGE esq_119_013
% BEGIN PAGE esq_119_014
\scanpage{119/14}
\mspage{12}

line is defined over the absolute base field $\mathbb{Q}$ ,
as are the allowed ramification points, the algebraic curves
obtained are defined not only over $\mathbb{C}$, but over the algebraic
closure $\overline{\mathbb{Q}}$ of $\mathbb{Q}$ in $\mathbb{C}$. As for the original map, it
is recovered on the algebraic curve as the inverse image of the real
segment $[0,1]$ (where $0$ is regarded as a vertex, and $1$ as
the midpoint of a ``folded edge'' with 1 as its center), which constitutes
the ``universal oriented 2-map'' on the Riemann sphere (*).
The points of the algebraic curve $X$ above 0, 1 and $\infty$
are precisely the vertices and the ``centers'' of the edges and
faces respectively of the map $(X,K)$, and the orders of the vertices
and faces are precisely the multiplicities of the zeros and
poles of the rational function (defined over $\overline{\mathbb{Q}}$) on $X$,
expressing its structural projection to $\mathbb{P}^1_{\mathbb{C}}$.

This discovery, which technically amounts to so little,
made a very strong impression on me, and represents
a decisive turning point in the course of my reflections, in particular
a shift in my center of interest in mathematics, which
suddenly became sharply localized. I do not believe that a
mathematical fact has ever struck me as much as this one, or had
a comparable psychological impact \notecall{copy1-2}{call-esq_119_014-2-1}{($^2$)}. This surely has to do with the
so familiar, nontechnical nature of the objects under consideration,
of which any child's drawing scribbled on a scrap of paper (provided
the drawing is all in one piece) gives a perfectly
explicit example. With such a drawing are associated subtle arithmetic
invariants, which will be completely upset as soon as
one adds another stroke. Since we are dealing here with spherical maps, necessarily

\noindent\rule{4cm}{0.4pt}

(*) There is an analogous description of finite unoriented maps,
possibly with boundary, in terms of \uline{real} algebraic curves,
more precisely of a covering of $\mathbb{P}^1_{\mathbb{R}}$ ramified only at 0, 1, $\infty$,
the surface with boundary associated with such a covering being $X(\mathbb{C})/\tau$, where $\tau$
is complex conjugation. The ``universal'' unoriented map is
here the disk, or upper hemisphere of the Riemann sphere, equipped,
as before, with the embedded 1-complex $K = [0,1]$.

\margin{14}

\finishscan
% END PAGE esq_119_014
% BEGIN PAGE esq_119_015
\scanpage{119/15}
\mspage{13}

giving rise to curves of genus zero (which therefore
provide no ``moduli''), one can say that the curve
in question is ``pinned down'' as soon as three of its points are fixed, for
example three vertices of the map, or more generally three
centers of facets (vertices, edges or faces) - from that point on, the structural
map $f : X \to \mathbb{P}^{1}_{\mathbb{C}}$ can be interpreted as a
well-determined rational function
\[ f(z) = P(z)/Q(z) \in \mathbb{C}(z) \]
the quotient of two well-determined relatively prime polynomials,
with $Q$ monic, satisfying algebraic conditions
that express, in particular, the fact that $f$ is
unramified outside the values $0$, $1$, $\infty$, and imply that the
coefficients of these polynomials are \uline{algebraic numbers}; thus
their zeros are algebraic numbers, representing respectively
the vertices and the centers of the faces of the map under
consideration.

Returning to the general case, since finite maps are interpreted
as coverings over $\bar{\mathbb{Q}}$ of an algebraic curve defined
over the prime field $\mathbb{Q}$ itself, it follows that the Galois
group $\mathrm{I}\!\Gamma$ of $\bar{\mathbb{Q}}$ over $\mathbb{Q}$ acts naturally on the category of these
maps. For example, the action of an automorphism
$\gamma \in \mathrm{I}\!\Gamma$ on a spherical map given by the rational function
above is obtained by applying $\gamma$ to the coefficients of the
polynomials $P$, $Q$. Here, then, is this mysterious group $\mathrm{I}\!\Gamma$ intervening as
a transforming agent on topological-combinatorial shapes of the
most elementary kind imaginable, prompting questions
such as: are these given oriented maps ``conjugate'',
or: what exactly are the conjugates of this given oriented
map ? (There are, visibly, only finitely many.)

I have dealt with a few concrete cases (for coverings
of low degrees) by various expedients; J. Malgoire has dealt with
a few others - I doubt that there is a uniform method of
answering them by computer calculations. My reflection very quickly

\margin{15}

\finishscan
% END PAGE esq_119_015
% BEGIN PAGE esq_119_016
\scanpage{119/16}
\mspage{14}

set off in a more conceptual direction, to come to
grasp the nature of this action of $\mathrm{I\!\Gamma}$. One immediately sees
that, roughly speaking, this action is expressed by a certain
``outer'' action of $\mathrm{I\!\Gamma}$ on the profinite compactification\tnote{TN-016-01}{``Profinite compactification'' renders ``compactifié profini'', the expression used here. The source used ``complétions profinies'' on p.~6, translated there as ``profinite completions''.}
of the oriented cartographic group $\underline{C}_2^{+}$ , and this action in turn
is deduced, by passage to the quotient, from the canonical outer
action of $\mathrm{I\!\Gamma}$ on the profinite fundamental group $\hat{\pi}_{0,3}$ of $(U_{0,3})_{\bar{\mathbb{Q}}}$ ,
where $U_{0,3}$ denotes the standard curve of genus $0$ over the prime field
$\mathbb{Q}$ with three points removed. This is how my attention
turned toward what I have since called ``\uline{anabelian}
\uline{algebraic geometry}'', whose starting point is precisely a study
(for the moment restricted to characteristic zero) of the action
of ``absolute'' Galois groups (\unsure{in particular}\tnote{TN-016-02}{The French alternatives are ``nctamment|notamment''. Both are represented by ``in particular''; the letterform doubt remains linked to the source.}
the groups $\mathrm{Gal}(\bar{K}/K)$,
where $K$ is an extension of finite type of the prime field) on
geometric (profinite) fundamental groups of algebraic varieties
(defined over $K$), and more particularly (breaking with
a deeply rooted tradition) fundamental groups that are
very far from abelian groups (and that I therefore
call ``\uline{anabelian}''). Among these groups, and very close to the group
$\hat{\pi}_{0,3}$ , is the profinite compactification of the modular group $S\ell(2,\mathbb{Z})$,
whose quotient by the center $\pm 1$ contains the preceding group as
the congruence subgroup mod 2, and can moreover also be interpreted
as an oriented ``cartographic'' group, namely the one
that classifies \uline{triangulated} oriented maps (i.e. those whose
faces are triangles or monogons).

Every finite oriented map gives rise to a smooth projective
algebraic curve defined over $\bar{\mathbb{Q}}$ , and the question then immediately
arises: which algebraic curves over $\bar{\mathbb{Q}}$ are obtained
in this way - might one obtain them all, who knows ?
In more learned terms, might it be true that every smooth projective
algebraic curve defined over a number field would occur
as a possible ``modular \unsure{curve}''\tnote{TN-016-03}{The French reading remains ``ccurbe|courbe''. The English keeps the doubt on ``curve'' and supplies both source candidates here.}
parametrizing elliptic curves
equipped with a suitable rigidification ? Such a
supposition looked so crazy that I was almost

\margin{16}

\finishscan
% END PAGE esq_119_016
% BEGIN PAGE esq_119_017
\scanpage{119/17}
\mspage{15}

embarrassed to submit it to the experts in the subject. Deligne, consulted,
did indeed find the supposition crazy, but had no counterexample
up his sleeve. Less than a year later, at the International
Congress in Helsinki, the Soviet mathematician Bielyi\tnote{TN-017-01}{The source's spelling ``Bielyi'' is retained.}
announced precisely this result, with a proof of disconcerting simplicity
occupying two short pages of a letter from Deligne - probably never
has a deep and bewildering result been proved
in so few lines !

In the form in which Bielyi states it, his result says essentially
that \uline{every algebraic curve defined over a number}
\uline{field can be obtained as a covering of the projective line ramified}
\uline{only at the points} 0, 1, $\infty$. This result seems to have passed
more or less unnoticed. Yet it appears to me to have considerable
scope. For me, its essential message was that \uline{there is a deep}
\uline{identity between the combinatorics of finite maps on the one hand}, \uline{and the}
\uline{geometry of algebraic curves defined over number fields}
\uline{on the other}. This deep result, together with the algebraic-geometric
interpretation of finite maps, opens the door onto a new,
unexplored world - and within reach of everyone, who passes by without seeing it.

It was only almost three years later, seeing that the vast
horizons opening up here decidedly stirred
nothing in any of my students, or even in any of the three or
four colleagues of the highest caliber to whom I had the opportunity
to speak about it in detail, that I made a first exploratory
journey into this ``new world'', from January to June 1981. This first venture
took material form in a packet of some 1 300 manuscript pages,
christened ``The Long March through Galois Theory''.\tnote{TN-017-02}{The French title is ``La Longue Marche à travers la théorie de Galois''.}
It is above all an effort to understand the relations between
``arithmetic'' Galois groups and ``geometric'' profinite fundamental
groups. Fairly quickly it turned toward formulating
the action of $\mathrm{Gal}(\bar{\mathbb{Q}}/\mathbb{Q})$ on $\hat{\pi}_{0,3}$ in calculational terms, and at a
later stage, on the slightly larger group $S\ell(2,\mathbb{Z})^{\wedge}$, which
gives rise to a more elegant and more effective formalism. It was also
in the course of this work (but developed in separate notes)

\margin{17}

\finishscan
% END PAGE esq_119_017
% BEGIN PAGE esq_119_018
\scanpage{119/18}
\mspage{16}

that the central theme of anabelian algebraic geometry appeared:
reconstructing certain varieties $X$ called ``anabelian''
over an absolute field $K$ from their mixed fundamental group,
an extension of $\mathrm{Gal}(\bar{K}/K)$ by $\pi_1(X_{\bar{K}})$; it was then that the
``fundamental conjecture of anabelian algebraic geometry'' emerged,
close to the conjectures of Mordell and Tate that Faltings
has just proved\notecall{copy1-3}{call-esq_119_018-3-1}{($^3$)}. It was also there that a first reflection on
Teichmüller groups and the first intuitions about the manifold
structure of the ``Teichmüller tower'' began - the open moduli
multiplicities $M_{g,\nu}$ appearing, moreover, as the first important
examples, in dimension $> 1$, of varieties (or rather,
multiplicities) that seem well deserving of the name ``anabelian''.
Toward the end of this period of reflection, it appeared to me
as a foundational reflection on a theory still in
limbo, for which the name ``\uline{Galois-Teichmüller theory}''
seemed more appropriate than ``Galois Theory'', the name I had initially
given to my notes.

This is not the place to give a more detailed account
of this whole collection of questions, intuitions, ideas - including
tangible results, certainly. The most important seems to me to be the one
mentioned in passing in par. 2, namely the faithfulness of the outer
action of $\mathrm{I}\!\Gamma = \mathrm{Gal}(\bar{\mathbb{Q}}/\mathbb{Q})$ (and of its open subgroups) on $\hat{\pi}_{0,3}$,
and more generally (if I remember correctly) on the fundamental
group of any ``anabelian'' algebraic curve (i.e. whose genus
$g$ and ``number of holes'' $\nu$ satisfy the inequality $2g + \nu \geq 3$,
i.e. such that $\chi(X) < 0$) defined over a finite extension of $\mathbb{Q}$. This
result can be regarded as essentially equivalent to
Bielyi's theorem - it is the first concrete manifestation, through
a precise mathematical statement, of the ``message'' discussed
above.

I would like to finish this rapid overview with a few words
of commentary on the truly unimaginable richness of a typical
anabelian group such as the group $S\ell(2,\mathbb{Z})$ - probably the most
remarkable infinite discrete group one has encountered, which appears

\margin{18}

\finishscan
% END PAGE esq_119_018
% BEGIN PAGE esq_119_019
\scanpage{119/19}
\mspage{17}

in a multitude of guises (some of which have been touched on
in the present report), and which, from the viewpoint of
Galois-Teichmüller theory, can be regarded as the fundamental
``building block'' of the ``Teichmüller tower''. The structural
element of $S\ell(2,\mathbb{Z})$ that fascinates me above all is of course
the outer action of the Galois group $\mathrm{I}\!\Gamma$ on its profinite
compactification. By Bielyi's theorem, taking the profinite
compactifications of subgroups of finite index of $S\ell(2,\mathbb{Z})$, and the induced
outer action (passing also to an open subgroup of $\mathrm{I}\!\Gamma$
as necessary), \uline{one essentially obtains the fundamental groups}
\uline{of all algebraic curves} (not necessarily compact)
\uline{defined over number fields} $K$, \uline{and the outer action of}
$\mathrm{Gal}(\bar{K}/K)$ \uline{on them} - at least it is true that every such fundamental
group appears as a quotient of one of the former groups (*).
Taking into account the ``anabelian yoga'' (which remains conjectural),
according to which an anabelian algebraic curve over a number field $K$
(a finite extension of $\mathbb{Q}$) is known up to isomorphism when
its mixed fundamental group is known (or, what amounts to the same thing,
the outer action of $\mathrm{Gal}(\bar{K}/K)$ on its geometric profinite
fundamental group), one can thus say that \uline{all algebraic curves}
\uline{defined over number fields are} ``\uline{contained}'' \uline{in the profinite}
\uline{compactification} $S\ell(2,\mathbb{Z})\hat{}$, \uline{and in the knowledge of a certain sub-}%
\uline{group} $\mathrm{I}\!\Gamma$ \uline{of its group of outer automorphisms} !
Passing to the abelianizations of the preceding fundamental groups, one sees,
in particular, that all the abelian $\ell$-adic representations dear
to Tate and his associates, defined by Jacobians and generalized
Jacobians of algebraic curves defined over number fields,
are contained in this single action of $\mathrm{I}\!\Gamma$ on the anabelian profinite
group $S\ell(2,\mathbb{Z})\hat{}$!\notecall{copy1-4}{call-esq_119_019-4-1}{($^{4}$)}

\rule{5cm}{0.4pt}

(*) In fact, these are quotients of a particularly trivial
kind, by abelian subgroups that are products of ``Tate modules''
$\hat{\mathbb{Z}}(1)$, corresponding to ``loop groups'' about points at infinity.

\margin{19}

\finishscan
% END PAGE esq_119_019
% BEGIN PAGE esq_119_020
\scanpage{119/20}
\mspage{18}

Faced with this, some are content to shrug their shoulders
with a disillusioned air and bet that nothing can be drawn
from it all except dreams. They forget, or are unaware,
that our science, and every science, would amount to very little
if, since its beginnings, it had not been nourished by the dreams and visions
of those who devote themselves to it with passion.

4.        From the beginning of my reflection on two-dimensional maps,
I was especially interested in the so-called
``\uline{regular}'' maps, that is, those whose automorphism group
acts transitively (and consequently, simply transitively)
on the set of frames. In the oriented case and in terms of
the algebraic-geometric interpretation in the preceding paragraph,
these are the maps that correspond to a \uline{Galois} covering of
the projective line. Very quickly, too, and even before the connection
with algebraic geometry appeared, it also became
necessary not to exclude infinite maps, which in particular enter
naturally as universal coverings of
finite maps. It appears (as an immediate consequence of the ``dictionary''
of maps, extended to maps that need not be finite)
that for every pair of natural numbers $p, q$ $\geq$ 1, there exists up to
(nonunique) isomorphism one and only one 1-\uline{connected} map
of type $(p,q)$, i.e. all of whose vertices have order $p$
and all of whose faces have order $q$, and this map is regular.
It is pinned down by the choice of a frame, and its
automorphism group is then canonically isomorphic to the quotient
of the cartographic group (resp. of the oriented cartographic
group in the oriented case) by the additional relations
\[ \rho_s^{\,p} = \rho_f^{\,q} = 1. \]
The case in which this group is finite is the ``Pythagorean'' case of
regular spherical maps; the case in which it is infinite gives the tessellations

\margin{20}

\finishscan
% END PAGE esq_119_020
\end{document}
